\documentclass[12pt,a4paper,oneside]{article}
\usepackage[brazil]{babel}
\usepackage[utf8]{inputenc}
\usepackage[dvipsnames]{xcolor}
\usepackage{listings}
\usepackage{wrapfig}
\usepackage{olymp}

\setcounter{problem}{1}

\begin{document}

\begin{problem}{Palavra por palavra}{palavras}{2}

\begingroup
\setlength{\intextsep}{0pt}%
%\setlength{\columnsep}{0pt}

Uma equipe está trabalhando na decoração de um evento e precisa imprimir vários cartazes com diferentes frases.
Os cartazes são verticais, então cada palavra da frase deve aparecer em uma linha.
Faça um programa para automatizar o trabalho da equipe.

% \Notes

\InputFile

A entrada contém apenas uma linha com a frase que deve ser processada. A frase possui no máximo $200$ caracteres e contém apenas letras e espaços, sendo terminada por um ponto final. Os espaços aparecem apenas entre palavras e não há espaços seguidos.

\endgroup

\OutputFile

Escreva a sequência de palavras da frase, uma por linha.

% \Notes

\Example

\begin{example}
\exmp{
Hello world.
}{
Hello
world
}%
\end{example}

\begin{example}
\exmp{
Ciencia da Computacao.
}{
Ciencia
da
Computacao
}%
\end{example}

\begin{example}
\exmp{
To be or not to be.
}{
To
be
or
not
to
be
}%
\end{example}

%Comentários sobre o exemplo, se necessário.

\end{problem}

\end{document}